\documentclass[12pt,a4paper]{article}
\usepackage[utf8]{inputenc}
\usepackage[T1]{fontenc}
\usepackage[slovene]{babel}
\usepackage{graphicx}
\usepackage{url}
\usepackage[margin=2.5cm]{geometry}

\graphicspath{{slike/}}

\newcommand{\slika}[2]{%
  \IfFileExists{slike/#1}
    {\includegraphics[width=#2\textwidth]{#1}}
    {\fbox{\parbox[c][6cm][c]{0.85\textwidth}{\centering Manjka slika:\\ \texttt{slike/#1}}}}}

\begin{document}

\begin{titlepage}
  \centering
  {\large Univerza v Mariboru\\
   Fakulteta za elektrotehniko, računalništvo in informatiko\\[0.3cm]
   Informatika in tehnologije komuniciranja\\}
  \vfill
  {\Huge\bfseries Spletna stran za Pizzerijo Bumerang\\[0.4cm]}
  {\Large Pregled obstoječih rešitev\\[0.3cm]}
  {\large Naloga 6: Priprava na razvoj\\}
  \vfill
  {\large
   Nejc Hribernik (vodja)\\
   Aleks Rožič\\
   Anej Sgerm\\}
  \vfill
  {\large Maribor, november 2026}
\end{titlepage}

\tableofcontents
\newpage

\section{Uvod}

Pizzerija Bumerang je pizzerija v Dravogradu. Naš cilj je zanjo izdelati
spletno stran, na kateri bodo obiskovalci lahko spoznali lokal, pregledali
jedilnik, oddali naročilo ter rezervirali mizo.

Preden začnemo z razvojem, smo pregledali tri obstoječe rešitve, ki ponujajo
podobne storitve: spletno stran mednarodne verige Papa John's, spletno stran
mariborske Pizzerie Toscana ter slovensko platformo za dostavo hrane
ehrana.si. Pri vsaki smo opisali, čemu je namenjena, katere so njene ključne
funkcionalnosti ter kakšne so njene prednosti in slabosti. Ugotovitve smo
uporabili pri oblikovanju seznama funkcionalnosti in žičnih modelov
(wireframov) naše spletne strani.

\section{Obstoječa rešitev 1: Papa John's (papajohns.de)}

\subsection*{Kratek opis}
Papa John's je mednarodna veriga pizzerij. Njena nemška spletna stran omogoča
naročanje pic in drugih jedi z dostavo ali osebnim prevzemom. Namenjena je
strankam, ki želijo hrano hitro naročiti prek spleta, brez klica
\cite{papajohns}.

\subsection*{Ključne funkcionalnosti}
\begin{itemize}
  \item vnos naslova (ulica, hišna številka, poštna številka) ali uporaba trenutne lokacije,
  \item izbira med dostavo (\textit{Delivery}) in prevzemom (\textit{Collection}),
  \item prikaz predvidenega časa dostave (od 30 minut),
  \item prikaz akcij in paketov (npr.\ pica in prigrizek ali sladica),
  \item pregled jedilnika po kategorijah (npr.\ toplo pečeni sendviči),
  \item menjava jezika med nemščino in angleščino.
\end{itemize}

\subsection*{Prednosti}
\begin{itemize}
  \item stran je pregledna,
  \item naročanje se začne že na vrhu domače strani z vnosom naslova,
  \item akcije so velike in opazne,
  \item izbira med dostavo in prevzemom je enostavna,
  \item na voljo je angleški jezik.
\end{itemize}

\subsection*{Slabosti}
\begin{itemize}
  \item velik oglasni pas zavzame veliko prostora,
  \item jedilnik je viden šele ob drsenju navzdol,
  \item nekatere akcije veljajo le v izbranih poslovalnicah.
\end{itemize}

\begin{figure}[h]
  \centering
  \slika{resitev1.png}{0.95}
  \caption{Domača stran Papa John's (papajohns.de) \cite{papajohns}.}
  \label{fig:papajohns}
\end{figure}

\section{Obstoječa rešitev 2: Pizzeria Toscana Maribor (toscana-mb.si)}

\subsection*{Kratek opis}
Pizzeria Toscana je samostojna pizzerija v središču Maribora, na Ruski
ulici~4, ki ponuja italijansko hrano s poudarkom na picah \cite{toscanaKupi}.
Turistični portal Visit Maribor jo uvršča med najboljše pizzerije v mestu
\cite{visitmaribor}. Lokal ima lastno spletno stran, ki je namenjena
predvsem predstavitvi lokala in posredovanju kontaktnih podatkov
\cite{toscana}.

\subsection*{Ključne funkcionalnosti}
\begin{itemize}
  \item predstavitev lokala in njegove ponudbe,
  \item pregled jedilnika,
  \item kontaktni podatki: naslov, telefonska številka in e-pošta,
  \item objavljen delovni čas,
  \item prikaz lokacije na zemljevidu,
  \item rezervacija mize in naročilo za s seboj po telefonu \cite{toscanaKupi}.
\end{itemize}

\subsection*{Prednosti}
\begin{itemize}
  \item lokal ima lastno spletno stran in s tem lastno predstavitev,
        ki je ne deli z drugimi restavracijami,
  \item uporabnik hitro najde osnovne informacije: kje je lokal, kdaj je
        odprt in kako ga kontaktira,
  \item stran gradi prepoznavnost lokala, ki ima dober sloves tudi na
        turističnih portalih \cite{visitmaribor},
  \item ni provizij platform za dostavo, saj stranka komunicira neposredno
        z lokalom.
\end{itemize}

\subsection*{Slabosti}
\begin{itemize}
  \item ni spletnega naročanja ali košarice, naročilo je mogoče le po
        telefonu \cite{toscanaKupi},
  \item rezervacija mize ni mogoča prek spleta, temveč samo s klicem
        \cite{toscanaKupi},
  \item stran je namenjena predvsem informiranju, zato uporabnik po ogledu
        jedilnika ne more takoj ukrepati,
  \item pri klicu v času največje gneče lahko stranka čaka ali ne
        dobi odgovora.
\end{itemize}

\begin{figure}[h]
  \centering
  \slika{resitev2.png}{0.95}
  \caption{Domača stran Pizzerie Toscana Maribor \cite{toscana}.}
  \label{fig:resitev2}
\end{figure}

\section{Obstoječa rešitev 3: ehrana.si}

\subsection*{Kratek opis}
ehrana.si je slovenska spletna platforma in mobilna aplikacija za naročanje
dostave hrane na dom. Uporabnik izbere restavracijo, sestavi naročilo in ga
prejme na izbrani naslov \cite{ehrana}. Storitev je bila leta 2024 ponovno
zagnana pod novim lastnikom, ki želi pokriti tudi manjše kraje in podeželje
ter podpirati lokalne gostince \cite{ehrana24ur}. Deluje v spletnem
brskalniku ter na napravah Android in iOS \cite{ehranaApp}.

\subsection*{Ključne funkcionalnosti}
\begin{itemize}
  \item vnos naslova dostave in prikaz restavracij, ki dostavljajo na ta naslov,
  \item brskanje po restavracijah glede na kraj ali vrsto kuhinje (pice, burgerji, azijska, veganska \dots),
  \item košarica in oddaja naročila,
  \item plačilo z gotovino, s kartico ali prek spleta,
  \item uporabniški račun s shranjenimi naslovi, priljubljenimi naročili in zgodovino naročil,
  \item prijava z e-pošto, Facebookom ali Googlom,
  \item stran s popusti in akcijami ter naročanje na študentske bone,
  \item obvestilo o predvidenem času dostave,
  \item komentarji in ocene uporabnikov.
\end{itemize}

\subsection*{Prednosti}
\begin{itemize}
  \item slovenska storitev v slovenskem jeziku z veliko izbiro restavracij po vsej državi,
  \item več načinov plačila, tudi gotovina ob dostavi,
  \item pregleden postopek v treh korakih: izberi restavracijo, izberi jed, spremljaj dostavo,
  \item shranjevanje priljubljenih naročil omogoča hitro ponovno naročanje.
\end{itemize}

\subsection*{Slabosti}
\begin{itemize}
  \item gre za splošno platformo, zato posamezna restavracija nima lastne
        predstavitve (zgodbe, fotografij lokala, posebnosti),
  \item ni možnosti rezervacije mize,
  \item spletna stran ne omenja prevzema v lokalu,
  \item ponudba je odvisna od tega, ali restavracija dostavlja na uporabnikov
        naslov, kar je v manjših krajih lahko omejeno.
\end{itemize}

\begin{figure}[h]
  \centering
  \slika{ehrana.png}{0.95}
  \caption{Domača stran spletne platforme ehrana.si \cite{ehrana}.}
  \label{fig:ehrana}
\end{figure}

\section{Zaključek}

Pregled je pokazal, da vse tri rešitve dobro podpirajo spletno naročanje
hrane ali predstavitev lokala, vendar nobena ne pokrije vsega, kar potrebuje
manjša lokalna pizzerija. Spletna stran Papa John's ima pregledno naročanje
ter izbiro med dostavo in prevzemom, a je namenjena veliki verigi in jo
prekrivajo oglasi. Pizzeria Toscana ima kot samostojen lokal lastno spletno
stran, vendar ta služi predvsem informiranju, saj sta naročilo in
rezervacija mize mogoča le po telefonu. Platforma ehrana.si ponuja veliko
izbiro in več načinov plačila, vendar restavracija na njej nima lastne
predstavitve in ne omogoča rezervacije mize.

Na naši spletni strani za Pizzerijo Bumerang bomo zato združili dobre lastnosti
pregledanih rešitev: lastno predstavitev lokala, pregleden jedilnik s
košarico, enostavno oddajo naročila ter obrazec za rezervacijo mize in
kontakt. Izogibali se bomo velikim oglasnim pasovom, ki zakrijejo vsebino,
jedilnik pa bo dostopen neposredno iz navigacije.

\begin{thebibliography}{9}
\bibitem{papajohns} Papa John's Deutschland. Dostopno na:
  \url{https://www.papajohns.de} (dostopano 22.~11.~2026).
\bibitem{toscana} Pizzeria Toscana Maribor. Dostopno na:
  \url{https://www.toscana-mb.si} (dostopano 22.~11.~2026).
\bibitem{toscanaKupi} Kupi.com: Pizzeria Toscana, Maribor. Dostopno na:
  \url{https://kupi.com/en-ae/explore/slovenia/maribor/pizzeria-toscana}
  (dostopano 22.~11.~2026).
\bibitem{visitmaribor} Visit Maribor: Best pizza is home in Maribor. Dostopno na:
  \url{https://www.visitmaribor.si/en/what-to-do/wine-and-culinary/restaurants/in-the-company-of-the-best/top-good-pizza/}
  (dostopano 22.~11.~2026).
\bibitem{ehrana} ehrana.si -- Dostava hrane na dom. Dostopno na:
  \url{https://www.ehrana.si/} (dostopano 22.~11.~2026).
\bibitem{ehranaApp} ehrana -- Dostava Hrane, App Store. Dostopno na:
  \url{https://apps.apple.com/app/id6504228246} (dostopano 22.~11.~2026).
\bibitem{ehrana24ur} 24ur.com: Ehrana je nazaj, pokriti želijo tudi
  podeželje, 19.~7.~2024. Dostopno na:
  \url{https://www.24ur.com/novice/gospodarstvo/ehrana-je-nazaj-pokriti-zelijo-tudi-podezelje.html}
  (dostopano 22.~11.~2026).
\end{thebibliography}

\end{document}